%% ==========================================================================
%%  5  Pareto optimality and unit subsidies
%%
%%  The Pareto subsection of Section 3 of the arXiv version
%%  (arxiv/sections/main_result.tex) and its Appendix D
%%  (arxiv/sections/appendix_additive_objective_po.tex). Wording follows the
%%  arXiv text, with punctuation brought to the house style.
%% ==========================================================================
\section{Pareto Optimality and Unit Subsidies}
\label{sec:pareto}

We next examine whether the unit-subsidy guarantee can be obtained together with
Pareto optimality of the underlying allocation. Pareto optimality is evaluated
with respect to the allocation of items, without taking subsidies into account.
Formally, a complete allocation $\allocB$ \emph{Pareto dominates} a complete
allocation $\alloc$ if
\[
  v_i(\allocB_i)\geq v_i(\bundle i)
  \qquad\text{for every }i\in\agents,
\]
with strict inequality for at least one agent. In cost form this reads
$\cost_i(\allocB_i)\leq \cost_i(\bundle i)$. A complete allocation is
\emph{Pareto optimal} if it is not Pareto dominated by any other complete
allocation. The quantifier ranges over complete allocations: for chores,
withholding an item weakly benefits every agent, so allowing incomplete
allocations would leave no allocation Pareto optimal.

% --------------------------------------------------------------------------
\subsection{Negative dichotomous valuations}
\label{sec:pareto-chores}

For the binary additive subclass, Pareto optimality is compatible with the
unit-subsidy guarantee. This parallels the corresponding phenomenon for binary
additive goods, where Halpern and Shah~\cite{HS19} obtain a Pareto-efficient
allocation requiring total subsidy at most $n-1$. See also Goko et
al.~\cite{goko2024fair} for unit subsidies together with utilitarian optimality
in the more general matroidal goods setting. For completeness, we record the
simple chore-side observation.

\begin{proposition}
\label{prop:additive-po}
For binary additive costs, there exists a complete allocation $\mathcal A$ and
$p\in\{0,1\}^n$ such that $(\mathcal A,p)$ is envy-free,
$\sum_{i\in N}p_i\le n-1$, and $\mathcal A$ is Pareto optimal and EF1. Moreover,
such an outcome can be computed in polynomial time.
\end{proposition}

\begin{proof}
Let
\[
  U:=\{e\in M:c_i(\{e\})=1\text{ for every }i\in N\}
\]
be the universally costly chores. Assign every $e\notin U$ to an agent for whom
$c_i(e)=0$. Write $|U|=qn+r$, where $0\le r<n$, and distribute the chores in $U$
so that $r$ agents receive $q+1$ of them and the remaining agents receive $q$.

Set $p_i=1$ for the former $r$ agents and $p_i=0$ for the others. Then
$c_i(A_i)-p_i=q$ for every $i$. Moreover, every bundle $A_j$ contains $q+p_j$
universally costly chores, and hence
\[
  c_i(A_j)-p_j\ge q=c_i(A_i)-p_i
\]
for all $i,j$. Thus $(\mathcal A,p)$ is envy-free and $\sum_i p_i=r\le n-1$.

Every chore outside $U$ is allocated at zero cost, whereas every chore in $U$
necessarily contributes one unit of cost under every allocation. Hence
$\mathcal A$ minimizes total cost and is therefore Pareto optimal. Finally, an
agent receiving $q$ universally costly chores is already envy-free, while an
agent receiving $q+1$ becomes envy-free after removing one universally costly
chore. Thus $\mathcal A$ is EF1. The construction clearly requires only
polynomially many value queries and operations.
\end{proof}

Thus, in the binary additive domain, Pareto optimality is compatible with the
optimal unit-subsidy guarantee. The total-subsidy bound remains tight even under
Pareto optimality: in the identical binary additive instance $c_i(S)=|S|$ of
Example~\ref{ex:lowerbound}, every complete allocation is Pareto optimal, while
every envy-free solution requires total subsidy at least $n-1$. This
compatibility, however, breaks down once nonadditivity is allowed, even for
identical binary submodular costs.

\begin{proposition}
\label{prop:po-impossibility}
There exist negative dichotomous instances, even with two agents and identical
binary submodular cost functions, for which no Pareto-optimal allocation admits
an envy-free subsidy vector in $\{0,1\}^n$.
\end{proposition}

\begin{proof}
Consider two agents $\agents=\set{1,2}$ and three chores $\items=\set{a,b,c}$.
Both agents have the same cost function
\[
  \cost_1(S)=\cost_2(S)=\min\{\abs S,2\}.
\]
Every marginal cost belongs to $\{0,1\}$, so the cost function is dichotomous.
Moreover, it is submodular, since $k\mapsto\min\{k,2\}$ is a nondecreasing
concave function of cardinality. It is not additive, since
\[
  \cost_i(\set{a,b,c})=2<3
  =\sum_{g\in\set{a,b,c}}\cost_i(\set g).
\]
Every complete allocation has, up to symmetry, either bundle-size profile
$(3,0)$ or $(2,1)$. First, consider the allocation in which agent $1$ receives
all three chores. Its cost profile is $(2,0)$. The possible cost profiles of
complete allocations are
\[
  (2,0),\qquad (2,1),\qquad (1,2),\qquad (0,2),
\]
and none of the latter three Pareto dominates $(2,0)$. Hence, the concentrated
allocation is Pareto optimal. By symmetry, the allocation in which agent $2$
receives all three chores is also Pareto optimal. However, a concentrated
allocation cannot be made envy-free with unit subsidies. If agent $1$ receives
all three chores, envy-freeness requires
\[
  2-p_1\leq -p_2,
\]
and hence
\[
  p_1-p_2\geq 2.
\]
Thus, no subsidy vector in $\{0,1\}^2$ suffices. The case in which agent $2$
receives all three chores is symmetric.

Now consider a split allocation, with bundle sizes $(2,1)$ up to symmetry. Its
cost profile is $(2,1)$, so it can be made envy-free with subsidies $(1,0)$, up
to symmetry. Nevertheless, it is not Pareto optimal. Transferring the singleton
chore to the agent holding two chores changes the cost profile from $(2,1)$ to
$(2,0)$: the cost of one agent remains unchanged, while the other agent's cost
strictly decreases.

Thus, the Pareto-optimal allocations are precisely the two concentrated
allocations, and each requires a subsidy of at least $2$ for the agent receiving
all three chores. Conversely, every allocation that admits an envy-free subsidy
vector in $\{0,1\}^2$ is Pareto dominated. Therefore, no allocation is
simultaneously Pareto optimal and envy-free with subsidies in $\{0,1\}^2$.
\end{proof}

All eight complete allocations of this instance, with their minimal subsidies,
are listed in Example~\ref{app:pareto} of Appendix~\ref{app:pareto-sec}. Thus,
the unit-subsidy guarantee of Theorem~\ref{thm:m-main} cannot, in general, be
strengthened to require Pareto optimality.

% --------------------------------------------------------------------------
\subsection{Additive objective signed dichotomous valuations}
\label{sec:pareto-objective}

For objective signed dichotomous valuations (Definition~\ref{def:objective-signed}),
Proposition~\ref{prop:po-impossibility} already rules out Pareto optimality in
general, since negative dichotomous valuations form a subclass. Under additivity,
however, Pareto optimality can be recovered together with the subsidy bound of
Theorem~\ref{thm:objective-mixed}.

\begin{proposition}
\label{prop:additive-objective-po}
For every instance with additive objective signed dichotomous valuations, there
exist a complete allocation $\alloc$ and a subsidy vector $p\in\{0,1\}^n$ such
that $(\alloc,p)$ is envy-free,
\[
  \sum_{i\in N}p_i\leq n-1,
\]
and $\alloc$ is Pareto optimal. Moreover, such an outcome can be computed in
polynomial time.
\end{proposition}

\begin{proof}
The proof has three steps. First, allocate the goods using Halpern and Shah's
binary-additive result, obtaining an envy-freeable allocation with unit
subsidies. Next, assign every chore that is costless to some agent to such an
agent. Finally, assign each universally costly chore to an agent currently
receiving zero subsidy and raise that agent's subsidy from $0$ to $1$. If all
subsidies become $1$, subtract one from every subsidy. This preserves
envy-freeness and keeps total subsidy at most $n-1$. Since every item is assigned
to an agent attaining its maximum possible value, the final allocation is Pareto
optimal.

We now formally write the details. Let $M=G\mathbin{\dot\cup}C$ be the common
partition into goods and chores. Since the valuations are additive, for every
agent $i$,
\[
  v_i(g)\in\{0,1\}\qquad (g\in G)
\]
and
\[
  v_i(c)\in\{0,-1\}\qquad (c\in C).
\]

We first allocate the goods. Ignore, for the moment, goods that have value zero
for every agent, and apply the binary-additive goods result of Halpern and
Shah~\cite{HS19} to the remaining goods. Their construction returns a
non-wasteful allocation $\alloc^G$ that is envy-freeable and whose minimal
subsidies have total at most $n-1$. Moreover, their proof shows that every path
in the corresponding weighted envy graph has weight at most one. Since all edge
weights are integral, the pointwise-minimal subsidies are integral as well, and
hence we may choose
\[
  p\in\{0,1\}^n,
  \qquad
  \sum_{i\in N}p_i\leq n-1.
\]
Goods having value zero for every agent may now be allocated arbitrarily,
without affecting either envy-freeness or the subsidy vector.

We next allocate the chores. Partition $C$ into
\[
  C^0:=\{c\in C:\text{there exists }i\in N\text{ with }v_i(c)=0\}
\]
and
\[
  C^-:=\{c\in C:v_i(c)=-1\text{ for every }i\in N\}.
\]
Thus, $C^-$ is the set of universally costly chores.

First, consider a chore $c\in C^0$. Choose an agent $k$ with $v_k(c)=0$ and
assign $c$ to $k$. This preserves envy-freeness with the same subsidy vector.
Indeed, agent $k$'s value for her own bundle is unchanged, whereas for every
agent $i\neq k$,
\[
  v_i(A_k\cup\{c\})
  =
  v_i(A_k)+v_i(c)
  \leq v_i(A_k).
\]
Hence, $k$'s bundle becomes weakly less desirable to every other agent, while all
other relevant bundle values remain unchanged. Repeating this operation allocates
all chores in $C^0$ while preserving the envy-free solution $(\alloc,p)$.

It remains to allocate the universally costly chores in $C^-$. We maintain the
following invariant:
\begin{equation}\label{eq:additive-mixed-invariant}
  (\alloc,p)\text{ is envy-free},\qquad
  p\in\{0,1\}^n,
  \qquad
  \text{and }p_k=0\text{ for some }k\in N.
\end{equation}
The invariant holds initially. In particular, the pointwise-minimal subsidy
vector has at least one zero component.

Let $c\in C^-$ be an unallocated chore, and choose an agent $k$ with $p_k=0$.
Assign $c$ to $k$ and define
\[
  \widehat p_k:=1,
  \qquad
  \widehat p_i:=p_i\quad(i\neq k).
\]
We claim that the resulting solution remains envy-free.

Since $c$ has value $-1$ for every agent,
\[
  v_k(A_k\cup\{c\})+\widehat p_k
  =
  v_k(A_k)-1+1
  =
  v_k(A_k)+p_k.
\]
Thus, agent $k$'s subsidized value for her own bundle is unchanged. Moreover,
for every $i\neq k$,
\[
  v_i(A_k\cup\{c\})+\widehat p_k
  =
  v_i(A_k)-1+1
  =
  v_i(A_k)+p_k.
\]
Hence, the subsidized value of $k$'s bundle is unchanged from the viewpoint of
every agent. All other bundles and subsidies are unchanged. Consequently, every
envy-freeness inequality is exactly the same as before, and the new solution is
envy-free.

If $\widehat p$ has a zero component, set $p:=\widehat p$ and continue.
Otherwise, if $\widehat p_i=1$ for all $i\in N$, replace it with
\[
  p_i:=\widehat p_i-1=0 \text{ for all } i\in N.
\]
Subtracting the same amount from every agent's subsidy leaves all envy-freeness
inequalities unchanged. Thus, the invariant \eqref{eq:additive-mixed-invariant}
is restored. Repeating this procedure allocates every chore in $C^-$ and
produces a complete allocation $\alloc$ with
\[
  p\in\{0,1\}^n
  \qquad\text{and}\qquad
  \sum_{i\in N}p_i\leq n-1.
\]

It remains to prove Pareto optimality. By construction, every good that is
valued at one by some agent is assigned to an agent who values it at one. A good
valued at zero by every agent may be assigned arbitrarily. Every chore in $C^0$
is assigned to an agent who values it at zero, while every chore in $C^-$ has
value $-1$ for every agent regardless of its recipient. Therefore, every item is
assigned to an agent attaining the maximum possible value for that item. By
additivity,
\[
  \sum_{i\in N}v_i(A_i)
  =
  \sum_{e\in M}\max_{i\in N}v_i(e),
\]
which is the maximum possible utilitarian welfare over all complete allocations.
Hence $\alloc$ is Pareto optimal.

Finally, the binary-additive goods allocation and its supporting subsidies can
be computed in polynomial time by the construction of Halpern and
Shah~\cite{HS19}. The subsequent chore-allocation steps require only identifying
a zero-valued recipient for each chore and updating a binary subsidy vector.
Hence, the entire construction runs in polynomial time.
\end{proof}
